% --- Patron: The Carrion-King (Decay, Renewal & Transformation) ---
\subsection{The Carrion-King (Decay, Renewal \& Transformation)}

\textit{Lore.} The Carrion-King is the master of endings that become beginnings. He does not destroy, but transforms—turning death into new life, decay into opportunity, and endings into fresh starts. His followers are harvesters of potential, seeing in every fall the seeds of future growth.

\begin{quote}
What crumbles feeds what grows. What dies becomes the soil of tomorrow's triumph.
\end{quote}

\subsubsection*{Rites of the Carrion-King}

\paragraph{Rite of Consuming Rot (Low, 5 XP)} \hfill \\
\emph{Instant; Touch; Yes (decay only).} \\
\textbf{Materials:} Organic matter in early stages of decay. \\
\textbf{Effect:} Accelerate natural decay to weaken or destroy: \textbf{+2 Effect} to \emph{Break/Sabotage} on organic materials (ropes, leather, wood). Gain 1 Boon if the decay creates an opportunity for you or allies. \\
\textbf{Push It:} Spread decay to similar materials in Close range; mark \textbf{1 SB (Clubs)} as the rot becomes noticeable. \\
\emph{Requires: Familiar (\textit{Invoke:} 1 Boon).}

\paragraph{Rite of the Harvested End (Low, 4 XP)} \hfill \\
\emph{Scene; Touch; No.} \\
\textbf{Materials:} The remains of a recently ended thing (burnt letter, wilted flower, shattered glass). \\
\textbf{Effect:} Extract value from endings: from a defeated enemy, gain \textbf{+1 die} to next action; from a failed plan, re-roll one 1 on your next roll; from a broken item, gain 1 SB to spend immediately. \\
\textbf{Push It:} Harvest additional value but mark \textbf{Fatigue 1} from dwelling on endings. \\
\emph{Requires: Familiar (\textit{Invoke:} 1 Boon).}

\paragraph{Rite of the Fertile Death (Standard, 8 XP)} \hfill \\
\emph{Scene; Zone; No.} \\
\textbf{Materials:} Ashes, compost, or the remains of anything that once lived. \\
\textbf{Effect:} Transform death into growth: create beneficial terrain (cover, concealment, or advantageous positioning) \textbf{OR} grant allies \textbf{+1 die} to healing/recovery rolls. Choose one effect per scene. \\
\textbf{Push It:} Both effects apply but attract unwanted attention (vermin, scavengers, or curious onlookers). \\
\emph{Requires: Familiar + Codex (\textit{Invoke:} 1 Boon).}

\paragraph{Rite of the Transformed Spirit (Standard, 7 XP)} \hfill \\
\emph{Instant; Near; No.} \\
\textbf{Materials:} A token from a deceased being (hair, nail, written name). \\
\textbf{Effect:} Channel the essence of what was: gain one skill die reflecting the deceased's expertise for one scene \textbf{OR} ask one question about their knowledge/abilities. \\
\textbf{Push It:} The spirit's influence lingers – gain permanent insight (\textbf{+1 die} specialty) but suffer occasional possession‑like effects (GM discretion). \\
\emph{Requires: Familiar + Codex (\textit{Invoke:} 1 Boon).}

\paragraph{Rite of the Great Consumption (High, 13 XP)} \hfill \\
\emph{Scene; Zone; No.} \\
\textbf{Materials:} A significant amount of organic matter (corpse, fallen tree, collapsed building). \\
\textbf{Effect:} Transform a large area through decay and renewal: choose two – create difficult terrain that favours you, summon Cap 3 swarm of scavengers as temporary allies, or generate valuable reagents worth 2 XP. \\
\textbf{Push It:} All three effects occur but start a 6‑segment \textbf{Ecosystem Disruption} clock that will cause problems later. \\
\emph{Requires: Familiar + Codex + Tier III (\textit{Invoke:} \textbf{2 Boons}).} \\
\emph{Obligation:} 7 segments.

\paragraph{Rite of the Eternal Cycle (High, 14 XP)} \hfill \\
\emph{Extended; Touch; No.} \\
\textbf{Materials:} The complete remains of something significant that has ended. \\
\textbf{Effect:} Complete a transformation cycle: destroy one major asset/enemy/obstacle and create something new of equal or greater value. GM and player collaborate to define the transformation. \\
\textbf{Push It:} The transformation is immediate and spectacular but creates a 6‑segment \textbf{Cycle Debt} clock – the King will demand another significant ending soon. \\
\emph{Requires: Familiar + Codex + Tier III (\textit{Invoke:} \textbf{2 Boons}).} \\
\emph{Obligation:} 7 segments.

\subsubsection*{Cantors \& Cults of the Carrion-King}

\textbf{The Cantor's Song.} A Cantor of the Carrion-King sings the \textit{dirge of transformation} – melodies that begin in mourning and end in fertile silence. Their voice carries the scent of turned earth and the weight of things laid to rest. They are called to funerals, to the edges of battlefields, and to compost heaps where the old year is buried. A Cantor of the King does not seek applause; they seek the moment when grief becomes acceptance, and acceptance becomes growth.

\textbf{The Cult of the Compost-Circle.} The King’s cults are not secret, but they are overlooked – gardeners, gravediggers, bone‑setters, and midwives who have seen too many births and deaths to fear the space between. Their rituals are practical: the annual turning of the village burial mound, the blessing of a new compost heap, the singing of a child into the world while its grandmother is sung out. The Compost-Circle keeps no written creed, only a shared understanding: \textit{“Nothing is wasted. Not tears, not bones, not the silence after a name is spoken.”}

\subsubsection*{Witchcraft of the Carrion-King (Threshold)}

A witch who walks with the Carrion-King is a \textit{rot‑warden}, a keeper of the slow alchemy that turns decay into soil. Their magic is patient, overlooked, and often performed in the half‑light of a root cellar or a fog‑bound cemetery.

\textbf{The Witch's Tool: The Bone‑Needle.} A needle carved from the long‑bone of a creature that died of old age. It can stitch wounds that will not heal, but more often it is used to “sew shut” a corpse’s mouth, ensuring the dead do not speak of things they saw in the dark. Once per session, touching the needle to a threshold (a door, a grave marker, a coffin lid) seals it against spirits for a scene.

\textbf{A Hedge Gift: Turn the Soil (2 XP).} Once per scene, you may touch a patch of fertile ground (or a pile of compost, a flowerpot, a grave) and immediately know which plant or creature will thrive there next season, or what was buried there last. The GM must answer one factual question about the site’s recent natural history.

\subsubsection*{Carrion-King’s Corruption Table}

\begin{tblr}{
  colspec = {Q[c,1cm] X[2.5] X[2.5]},
  rowhead = 1,
  row{odd} = {bg=gray!10},
  row{1} = {bg=gray!30, font=\bfseries},
  hlines,
}
Tier & Benefit & Cost / Quirk \\
1 & \textbf{Carrion's Insight:} +1 die to Notice decay or hidden weaknesses in structures or beings. & Must inspect decay first‑hand; suffer 1 Fatigue when exposed to fresh death or rot. \\
2 & \textbf{Deathward Sense:} Once per session, detect the last living moment of a dead being within Close range. & Cannot lie about death you’ve witnessed; must correct falsehoods. \\
3 & \textbf{Rotblood Resilience:} +1 die to resist disease and poison. & Immune system adapts slowly; each new disease/poison requires 1 Fatigue to resist. \\
4 & \textbf{Glean from Grief:} Once per scene, gain +1 die after witnessing a significant loss or defeat. & Compelled to linger at scenes of death; must spend one beat observing or risk 1 SB (Clubs). \\
5 & \textbf{Cycle's Whisper:} You can sense the “next ending” in any process – ask the Keeper one question about how a situation will collapse or conclude. & Must speak the truth about what you see, even if it harms your position. \\
6+ & \textbf{Eternal Bloom:} Once per session, declare a “death that births life.” Sacrifice an asset or ally to create something new of equal or greater value. & Mark +2 Obligation when using this power. \\
\end{tblr}